\clearpage

\section{SUBPART E --- DRUGS INTENDED TO TREAT LIFE-THREATENING AND SEVERELY-DEBILITATING ILLNESSES}

\textbf{Authority:} 21 U.S.C. 351, 352, 353, 355, 371; 42 U.S.C. 262.

\textbf{Source:} 53 FR 41523, Oct. 21, 1988, unless otherwise noted.

\subsection{\S~312.80 Purpose}

The purpose of this section is to establish procedures designed to expedite the development, evaluation, and marketing of new therapies intended to treat persons with life-threatening and severely-debilitating illnesses, especially where no satisfactory alternative therapy exists. As stated \S~314.105(c) of this chapter, while the statutory standards of safety and effectiveness apply to all drugs, the many kinds of drugs that are subject to them, and the wide range of uses for those drugs, demand flexibility in applying the standards. The Food and Drug Administration (FDA) has determined that it is appropriate to exercise the broadest flexibility in applying the statutory standards, while preserving appropriate guarantees for safety and effectiveness. These procedures reflect the recognition that physicians and patients are generally willing to accept greater risks or side effects from products that treat life-threatening and severely-debilitating illnesses, than they would accept from products that treat less serious illnesses. These procedures also reflect the recognition that the benefits of the drug need to be evaluated in light of the severity of the disease being treated. The procedure outlined in this section should be interpreted consistent with that purpose.

\subsubsection{Physical AI Adaptation of \S~312.80}

(a) The purpose of this section extends to the expedited development, evaluation, and marketing of Physical AI-assisted therapies intended to treat persons with life-threatening and severely-debilitating illnesses. The flexibility in applying statutory standards recognized in this section applies equally to Physical AI-assisted drug administration, provided that the Physical AI system meets the safety requirements of this part and has achieved the applicable USL threshold. The recognition that physicians and patients are generally willing to accept greater risks from products treating life-threatening illnesses extends to the Physical AI system risks associated with such treatments, including robotic system risks, AI/ML model uncertainty, and simulation-validated procedure risks.

[53 FR 41523, Oct. 21, 1988. Physical AI adaptation added 17 March 2026.]

\subsection{\S~312.81 Scope}

This section applies to new drug and biological products that are being studied for their safety and effectiveness in treating life-threatening or severely-debilitating diseases.

(a) For purposes of this section, the term ``life-threatening'' means:

\begin{description}[leftmargin=0.5cm,style=sameline]
\item[(1)] Diseases or conditions where the likelihood of death is high unless the course of the disease is interrupted; and
\item[(2)] Diseases or conditions with potentially fatal outcomes, where the end point of clinical trial analysis is survival.
\end{description}

(b) For purposes of this section, the term ``severely debilitating'' means diseases or conditions that cause major irreversible morbidity.

(c) Sponsors are encouraged to consult with FDA on the applicability of these procedures to specific products.

[53 FR 41523, Oct. 21, 1988, as amended at 64 FR 401, Jan. 5, 1999]

\subsection{\S~312.82 Early Consultation}

For products intended to treat life-threatening or severely-debilitating illnesses, sponsors may request to meet with FDA-reviewing officials early in the drug development process to review and reach agreement on the design of necessary preclinical and clinical studies. Where appropriate, FDA will invite to such meetings one or more outside expert scientific consultants or advisory committee members. To the extent FDA resources permit, agency reviewing officials will honor requests for such meetings.

(a) \textit{Pre-investigational new drug (IND) meetings.} Prior to the submission of the initial IND, the sponsor may request a meeting with FDA-reviewing officials. The primary purpose of this meeting is to review and reach agreement on the design of animal studies needed to initiate human testing. The meeting may also provide an opportunity for discussing the scope and design of phase 1 testing, plans for studying the drug product in pediatric populations, and the best approach for presentation and formatting of data in the IND.

(b) \textit{End-of-phase 1 meetings.} When data from phase 1 clinical testing are available, the sponsor may again request a meeting with FDA-reviewing officials. The primary purpose of this meeting is to review and reach agreement on the design of phase 2 controlled clinical trials, with the goal that such testing will be adequate to provide sufficient data on the drug's safety and effectiveness to support a decision on its approvability for marketing, and to discuss the need for, as well as the design and timing of, studies of the drug in pediatric patients. For drugs for life-threatening diseases, FDA will provide its best judgment, at that time, whether pediatric studies will be required and whether their submission will be deferred until after approval. The procedures outlined in \S~312.47(b)(1) with respect to end-of-phase 2 conferences, including documentation of agreements reached, would also be used for end-of-phase 1 meetings.

\subsubsection{Physical AI Adaptation of \S~312.82}

(c) \textit{Physical AI early consultation.} Early consultation meetings under this section for products that will involve Physical AI-assisted administration may additionally address:

\begin{description}[leftmargin=0.8cm,style=sameline]
\item[(i)] The design and scope of simulation validation studies for the Physical AI system, including the selection of simulation frameworks and the criteria for cross-framework behavioral consistency;
\item[(ii)] The proposed Physical AI system classification under \S~312.401 and the associated class-specific requirements;
\item[(iii)] The PCCP approach for managing anticipated AI/ML model updates during the investigation;
\item[(iv)] The proposed human oversight protocol, including the operator-to-system ratio and emergency stop specifications; and
\item[(v)] The approach for quantifying and managing the sim-to-real gap for the specific clinical procedures to be performed.
\end{description}

[53 FR 41523, Oct. 21, 1988, as amended at 63 FR 66669, Dec. 2, 1998. Physical AI adaptation added 17 March 2026.]

\subsection{\S~312.83 Treatment Protocols}

If the preliminary analysis of phase 2 test results appears promising, FDA may ask the sponsor to submit a treatment protocol to be reviewed under the procedures and criteria listed in \S\S~312.305 and 312.320. Such a treatment protocol, if requested and granted, would normally remain in effect while the complete data necessary for a marketing application are being assembled by the sponsor and reviewed by FDA (unless grounds exist for clinical hold of ongoing protocols, as provided in \S~312.42(b)(3)(ii)).

\subsubsection{Physical AI Adaptation of \S~312.83}

(a) \textit{Physical AI treatment protocols.} When a treatment protocol under this section involves Physical AI-assisted drug administration, the treatment protocol submission shall include a summary of the Physical AI system's performance data from the phase 2 investigation, the current USL score, and a plan for maintaining the Physical AI system at treatment sites with the same safety protections as used in the clinical investigation. The Physical AI system used under a treatment protocol must have achieved at least the USL threshold specified for the applicable procedure type.

[53 FR 41523, Oct. 21, 1988, as amended at 76 FR 13880, Mar. 15, 2011. Physical AI adaptation added 17 March 2026.]

\subsection{\S~312.84 Risk-Benefit Analysis in Review of Marketing Applications for Drugs to Treat Life-Threatening and Severely-Debilitating Illnesses}

(a) FDA's application of the statutory standards for marketing approval shall recognize the need for a medical risk-benefit judgment in making the final decision on approvability. As part of this evaluation, consistent with the statement of purpose in \S~312.80, FDA will consider whether the benefits of the drug outweigh the known and potential risks of the drug and the need to answer remaining questions about risks and benefits of the drug, taking into consideration the severity of the disease and the absence of satisfactory alternative therapy.

(b) In making decisions on whether to grant marketing approval for products that have been the subject of an end-of-phase 1 meeting under \S~312.82, FDA will usually seek the advice of outside expert scientific consultants or advisory committees. Upon the filing of such a marketing application under \S~314.101 or part 601 of this chapter, FDA will notify the members of the relevant standing advisory committee of the application's filing and its availability for review.

(c) If FDA concludes that the data presented are not sufficient for marketing approval, FDA will issue a complete response letter under \S~314.110 of this chapter or the biological product licensing procedures. Such letter, in describing the deficiencies in the application, will address why the results of the research design agreed to under \S~312.82, or in subsequent meetings, have not provided sufficient evidence for marketing approval. Such letter will also describe any recommendations made by the advisory committee regarding the application.

(d) Marketing applications submitted under the procedures contained in this section will be subject to the requirements and procedures contained in part 314 or part 600 of this chapter, as well as those in this subpart.

\subsubsection{Physical AI Adaptation of \S~312.84}

(e) \textit{Physical AI risk-benefit analysis.} When a marketing application involves a drug that was investigated using Physical AI-assisted administration, the risk-benefit analysis shall additionally consider:

\begin{description}[leftmargin=0.8cm,style=sameline]
\item[(i)] The Physical AI system's safety record throughout the clinical investigation, including all Physical AI adverse events reported under \S~312.32(g) and the system's USL score trajectory;
\item[(ii)] Whether the Physical AI-assisted administration provided a clinically meaningful improvement in drug delivery precision, treatment outcomes, or patient safety compared to conventional administration methods;
\item[(iii)] The robustness of the simulation validation data and the magnitude of any sim-to-real gap observed during the investigation; and
\item[(iv)] The availability and adequacy of the PCCP for managing post-market AI/ML model updates while maintaining the safety and effectiveness demonstrated during the investigation.
\end{description}

[53 FR 41523, Oct. 21, 1988, as amended at 73 FR 39607, July 10, 2008. Physical AI adaptation added 17 March 2026.]

\subsection{\S~312.85 Phase 4 Studies}

Concurrent with marketing approval, FDA may seek agreement from the sponsor to conduct certain postmarketing (phase 4) studies to delineate additional information about the drug's risks, benefits, and optimal use. These studies could include, but would not be limited to, studying different doses or schedules of administration than were used in phase 2 studies, use of the drug in other patient populations or other stages of the disease, or use of the drug over a longer period of time.

\subsubsection{Physical AI Adaptation of \S~312.85}

(a) \textit{Physical AI phase 4 studies.} Phase 4 studies for drugs approved with Physical AI-assisted administration may additionally include studies to evaluate: (i) the Physical AI system's long-term performance stability and any degradation in drug delivery precision over extended use; (ii) the Physical AI system's performance across a broader range of patient populations, including patients with anatomical variations, comorbidities, or other factors that may affect Physical AI system performance; (iii) the effectiveness of the PCCP in maintaining system safety and effectiveness through AI/ML model updates; and (iv) real-world Physical AI adverse event rates compared to those observed in the controlled clinical investigation.

[53 FR 41523, Oct. 21, 1988. Physical AI adaptation added 17 March 2026.]

\subsection{\S~312.86 Focused FDA Regulatory Research}

At the discretion of the agency, FDA may undertake focused regulatory research on critical rate-limiting aspects of the preclinical, chemical/manufacturing, and clinical phases of drug development and evaluation. When initiated, FDA will undertake such research efforts as a means for meeting a public health need in facilitating the development of therapies to treat life-threatening or severely debilitating illnesses.

\subsection{\S~312.87 Active Monitoring of Conduct and Evaluation of Clinical Trials}

For drugs covered under this section, the Commissioner and other agency officials will monitor the progress of the conduct and evaluation of clinical trials and be involved in facilitating their appropriate progress.

\subsubsection{Physical AI Adaptation of \S~312.87}

(a) \textit{Physical AI active monitoring.} For Physical AI-assisted drug investigations covered under this section, active monitoring by the Commissioner and agency officials may additionally include monitoring of the Physical AI system's clinical performance metrics, USL score changes, and simulation-to-clinical correlation data. FDA may request interim Physical AI performance data beyond the standard reporting requirements of \S~312.33 when such data are necessary to facilitate the appropriate progress of the investigation.

[53 FR 41523, Oct. 21, 1988. Physical AI adaptation added 17 March 2026.]

\subsection{\S~312.88 Safeguards for Patient Safety}

All of the safeguards incorporated within parts 50, 56, 312, 314, and 600 of this chapter designed to ensure the safety of clinical testing and the safety of products following marketing approval apply to drugs covered by this section. This includes the requirements for informed consent (part 50 of this chapter) and institutional review boards (part 56 of this chapter). These safeguards further include the review of animal studies prior to initial human testing (\S~312.23), and the monitoring of adverse drug experiences through the requirements of IND safety reports (\S~312.32), safety update reports during agency review of a marketing application (\S~314.50 of this chapter), and postmarketing adverse reaction reporting (\S~314.80 of this chapter).

\subsubsection{Physical AI Adaptation of \S~312.88}

(a) \textit{Physical AI patient safety safeguards.} In addition to the safeguards described in this section, the following Physical AI-specific safeguards apply to drugs covered by this section when Physical AI-assisted administration is used:

\begin{description}[leftmargin=0.8cm,style=sameline]
\item[(i)] The Physical AI system safety requirements of \S~312.60(b) through (g), including human oversight, pre-procedure verification, and emergency stop capabilities;
\item[(ii)] Physical AI adverse event reporting under \S~312.32(g);
\item[(iii)] The Physical AI System Description requirements of \S~312.23(g), including simulation validation and cybersecurity assessments;
\item[(iv)] The Physical AI record retention requirements of \S~312.57(f) and (g); and
\item[(v)] The Physical AI system classification and class-specific requirements of \S~312.401.
\end{description}

(b) The heightened flexibility recognized under \S~312.80 for drugs treating life-threatening and severely-debilitating illnesses does not diminish the Physical AI system safety requirements of this part. Regardless of the severity of the disease being treated, all Physical AI systems used in clinical investigations must meet the applicable USL threshold and must comply with the human oversight requirements of \S~312.404.

[53 FR 41523, Oct. 21, 1988. Physical AI adaptation added 17 March 2026.]

\clearpage

\section{SUBPART F --- MISCELLANEOUS}

\subsection{\S~312.110 Import and Export of Investigational New Drugs}

(a) \textit{Imports.} An investigational new drug offered for import into the United States complies with the requirements of this part if it is subject to an IND that is in effect for it under \S~312.40 and: (1) the consignee in the United States is the sponsor of the IND; (2) the consignee is a qualified investigator named in the IND; or (3) the consignee is the domestic agent of a foreign sponsor, is responsible for the control and distribution of the investigational drug, and the party who distributes the drug to the investigator is an authorized participant in the investigation. The importation of an investigational new drug or drug that is a biological product may be authorized by the Director of the Center for Drug Evaluation and Research or the Director of the Center for Biologics Evaluation and Research to permit the shipment of the drug for a clinical investigation authorized under an IND.

(b) \textit{Exports.} An investigational new drug may be exported from the United States for use in a clinical investigation under any of the following conditions:

\begin{description}[leftmargin=0.5cm,style=sameline]
\item[(1)] An IND is in effect for the drug under \S~312.40 and each person who receives the drug is an investigator named in the application; or the drug is being exported to further a clinical trial in one of the countries listed in section 802(b)(1)(A) of the act.
\item[(2)] The drug is being exported to a country not listed in section 802(b)(1)(A) of the act, provided that FDA determines that the drug is being investigated under an IND in the United States, that the drug complies with the laws of the importing country, and that the drug has valid marketing authorization in any country listed under section 802(b)(1)(A) of the act, as applicable.
\item[(3)] \textit{Expanded access for treatment use.} A drug that is not being commercially marketed in the United States may be exported if FDA determines that there is a clinical need, and certain conditions in the regulation are met.
\end{description}

\subsubsection{Physical AI Adaptation of \S~312.110}

(c) \textit{Import of Physical AI systems.} When a clinical investigation under an IND involves Physical AI systems that are manufactured or assembled outside the United States, the following additional requirements apply to the importation of such systems:

\begin{description}[leftmargin=0.8cm,style=sameline]
\item[(i)] The Physical AI system components shall be identified in the IND and shall comply with all applicable import requirements, including customs declarations that accurately describe the system components and their intended investigational use;
\item[(ii)] The Physical AI System Description shall identify the country of origin for all major hardware components, the location(s) where AI/ML models were trained, and the jurisdiction(s) whose data protection laws governed the training data;
\item[(iii)] The cybersecurity assessment submitted under \S~312.23(g)(iv) shall address supply chain security risks associated with imported Physical AI components, including firmware integrity verification and hardware tamper detection; and
\item[(iv)] The sponsor shall ensure that imported Physical AI systems comply with all applicable U.S. cybersecurity requirements and export control regulations, including restrictions on the re-export of encryption technology.
\end{description}

(d) \textit{Export of Physical AI systems.} When a Physical AI system is exported from the United States in connection with a clinical investigation:

\begin{description}[leftmargin=0.8cm,style=sameline]
\item[(i)] The export shall comply with all applicable export control regulations, including the Export Administration Regulations (15 CFR parts 730--774) as they apply to robotic systems, AI/ML technology, and encryption components;
\item[(ii)] The sponsor shall ensure that the Physical AI system will be used in the foreign clinical investigation in a manner consistent with the safety protections required under this part, including human oversight requirements and emergency stop capabilities; and
\item[(iii)] The sponsor shall maintain records of all Physical AI system exports, including the destination country, receiving investigator, system identifiers, and the date of export.
\end{description}

[52 FR 8831, Mar. 19, 1987, as amended at 52 FR 23031, June 17, 1987; 64 FR 401, Jan. 5, 1999; 67 FR 9586, Mar. 4, 2002. Physical AI adaptation added 17 March 2026.]

\subsection{\S~312.120 Foreign Clinical Studies Not Conducted Under an IND}

(a) \textit{Acceptance of studies.} FDA will accept as support for an IND or an application for marketing approval (an application under section 505(b) of the act or section 351 of the Public Health Service Act) a well-designed and well-conducted foreign clinical study not conducted under an IND, if the following conditions are met:

\begin{description}[leftmargin=0.5cm,style=sameline]
\item[(1)] The study was conducted in accordance with good clinical practice (GCP), including review and approval (or provision of a favorable opinion) by an independent ethics committee (IEC) that meets the requirements of the regulation.
\item[(2)] FDA is able to validate the data from the study through an onsite inspection if the agency deems it necessary.
\end{description}

(b) \textit{Supporting information.} A sponsor or applicant who wishes to rely on a foreign clinical study to support an IND or marketing application shall submit to FDA, among other things: (1) a description of the investigator's qualifications; (2) a description of the research facilities; (3) a detailed summary of the protocol and results of the study, and if the study is intended to support the effectiveness of a drug, information showing that the study is adequate and well controlled; (4) a description of the drug substance and formulation used in the study; (5) if the study is intended to support the safety of the drug, information about how investigators were informed about the safety of the drug, including adverse reaction information; (6) a statement that the study conformed to applicable GCP requirements; (7) the name and address of the IEC; and (8) a summary of the IEC's decision regarding the study.

(c) \textit{Waivers.} Requests for waivers of the requirements of paragraph (a)(1) shall be submitted in accordance with \S~312.10.

\subsubsection{Physical AI Adaptation of \S~312.120}

(d) \textit{Foreign Physical AI studies.} When a foreign clinical study not conducted under an IND involves Physical AI systems and is submitted to support an IND or marketing application, the following additional information shall be provided:

\begin{description}[leftmargin=0.8cm,style=sameline]
\item[(i)] A description of the Physical AI system used in the foreign study, including the system model, manufacturer, software version, AI/ML model version, and the Physical AI system's regulatory status in the country where the study was conducted;
\item[(ii)] Documentation that the Physical AI system used in the foreign study meets or is comparable to the safety requirements of this part, including human oversight protocols, emergency stop capabilities, and cybersecurity protections;
\item[(iii)] A USL assessment of the Physical AI system, conducted using the USL scoring methodology referenced in \S~312.1(e), or a demonstration that the foreign regulatory framework applied equivalent safety assessment criteria;
\item[(iv)] Simulation validation data for the Physical AI system, or a description of the pre-clinical validation methodology used and its comparability to the simulation validation requirements of \S~312.21(d);
\item[(v)] A description of any differences between the Physical AI system configuration used in the foreign study and the configuration proposed for use in the U.S. investigation, including hardware, software, AI/ML model, and operational protocol differences; and
\item[(vi)] Information sufficient to demonstrate that the Physical AI system data from the foreign study are reliable and that the Physical AI system performed consistently throughout the study.
\end{description}

(e) \textit{FDA validation of Physical AI data.} When FDA determines that onsite inspection is necessary under paragraph (a)(2), the inspection may include examination of the Physical AI system used in the foreign study, review of Physical AI telemetry data and system logs, and assessment of the Physical AI system's current operational status and compliance with the safety requirements described in the submission.

[52 FR 8831, Mar. 19, 1987, as amended at 52 FR 23031, June 17, 1987; 67 FR 9586, Mar. 4, 2002; 73 FR 22815, Apr. 28, 2008. Physical AI adaptation added 17 March 2026.]

\subsection{\S~312.130 Availability for Public Disclosure of Data and Information in an IND}

(a) The existence of an IND may be disclosed by FDA only to the extent that the agency determines that such disclosure would serve an important public purpose.

(b) After an IND is effective, FDA may disclose information about the IND as follows:

\begin{description}[leftmargin=0.5cm,style=sameline]
\item[(1)] If the existence of the IND has not been publicly disclosed or acknowledged, no data or information in the IND file are available for public disclosure.
\item[(2)] If the existence of the IND has been publicly disclosed or acknowledged, the following data and information in the IND file are immediately available for public disclosure upon request: (i) the name of the drug; (ii) the dosage form; (iii) the route of administration; (iv) a description of the drug's pharmacological and toxicological effects, if any, that are relevant to an assessment of whether the drug may be lawfully transported under section 505(i) of the act; and (v) a list of the active ingredients, if previously disclosed.
\item[(3)] After an IND has been terminated by the sponsor and all associated INDs have been terminated for 5 years, FDA may disclose additional information as specified in the regulation.
\end{description}

(c) Clinical investigators named on the application, or identified in an IND safety report, are not protected from disclosure by the agency.

(d) If an IND is discontinued and no marketing application is filed, data and information in the IND file are not available for public disclosure.

\subsubsection{Physical AI Adaptation of \S~312.130}

(e) \textit{Physical AI information disclosure.} The following provisions apply to Physical AI-specific data and information in an IND file:

\begin{description}[leftmargin=0.8cm,style=sameline]
\item[(i)] The Physical AI System Description, simulation validation data, and cybersecurity assessment submitted under \S~312.23(g) shall be treated as confidential commercial information and shall not be disclosed to the public without the sponsor's consent, except as otherwise provided by law;
\item[(ii)] When the existence of an IND involving a Physical AI system has been publicly disclosed, FDA may disclose that the investigation involves a Physical AI system and may identify the general type of Physical AI system (e.g., robotic surgical system, automated drug delivery system) without disclosing proprietary details of the system's design, algorithms, or cybersecurity architecture;
\item[(iii)] Aggregate Physical AI adverse event data may be disclosed in the same manner as other safety data under this section; and
\item[(iv)] USL scores and simulation validation results are confidential commercial information and shall not be disclosed without the sponsor's consent, except as required for public safety purposes.
\end{description}

[52 FR 8831, Mar. 19, 1987, as amended at 52 FR 23031, June 17, 1987; 67 FR 9586, Mar. 4, 2002. Physical AI adaptation added 17 March 2026.]

\subsection{\S~312.140 Address for Correspondence}

(a) Except as provided in paragraph (b) of this section, a sponsor shall send an initial IND submission to the Center for Drug Evaluation and Research or the Center for Biologics Evaluation and Research, as appropriate, at the addresses and in the manner described in \S~312.23(d).

(b) After FDA has assigned an IND number, the sponsor shall direct all subsequent IND submissions, amendments, reports, and other correspondence to the appropriate review division, identified by the IND number.

\subsubsection{Physical AI Adaptation of \S~312.140}

(c) \textit{Physical AI submissions.} All Physical AI-specific IND submissions, including the Physical AI System Description, Physical AI amendments, Physical AI adverse event reports, and Physical AI annual report supplements, shall be directed to the same review division as the parent IND. When submitting Physical AI-specific documents, the sponsor shall clearly identify the submission as Physical AI-related in the cover letter and shall reference the IND number and the specific Physical AI system(s) to which the submission pertains.

[52 FR 8831, Mar. 19, 1987, as amended at 52 FR 23031, June 17, 1987; 67 FR 9586, Mar. 4, 2002; 74 FR 13113, Mar. 26, 2009. Physical AI adaptation added 17 March 2026.]

\subsection{\S~312.145 Guidance Documents}

(a) FDA has made available guidance documents under \S~10.115 of this chapter to help sponsors and investigators comply with the requirements of this part.

(b) Guidance documents represent the agency's current thinking on a topic and do not create or confer any rights for or on any person and do not operate to bind FDA or the public.

\subsubsection{Physical AI Adaptation of \S~312.145}

(c) \textit{Physical AI guidance documents.} FDA may issue guidance documents specific to the use of Physical AI systems in clinical investigations, including guidance on:

\begin{description}[leftmargin=0.8cm,style=sameline]
\item[(i)] Preparation of the Physical AI System Description;
\item[(ii)] Simulation validation methodologies and acceptance criteria;
\item[(iii)] USL scoring methodology and threshold determination;
\item[(iv)] Physical AI adverse event classification and reporting;
\item[(v)] Cybersecurity requirements for Physical AI systems in clinical investigations;
\item[(vi)] Human oversight protocols for Physical AI-assisted drug administration;
\item[(vii)] AI/ML model change management under the PCCP framework; and
\item[(viii)] Digital twin validation for treatment planning applications.
\end{description}

Such guidance documents shall be developed with public input through the procedures described in \S~10.115 of this chapter and should consider relevant standards from recognized standards-developing organizations, including IEEE, ISO, and IEC.

[67 FR 9586, Mar. 4, 2002. Physical AI adaptation added 17 March 2026.]

\clearpage

\section{SUBPART G --- DRUGS FOR INVESTIGATIONAL USE IN LABORATORY RESEARCH ANIMALS OR IN VITRO TESTS}

\subsection{\S~312.160 Drugs for Investigational Use in Laboratory Research Animals or In Vitro Tests}

(a) \textit{Authorization to ship.}

\begin{description}[leftmargin=0.5cm,style=sameline]
\item[(1)] \hfill
\begin{description}[leftmargin=0.8cm,style=sameline]
\item[(i)] A person may ship a drug intended solely for tests in vitro or in animals used only for laboratory research purposes if it is labeled as follows: ``CAUTION: Contains a new drug for investigational use only in laboratory research animals, or for tests in vitro. Not for use in humans.''
\item[(ii)] A person may ship a biological product for investigational in vitro diagnostic use that is listed in \S~312.2(b)(2)(ii) if it is labeled as follows: ``CAUTION: Contains a biological product for investigational in vitro diagnostic tests only.''
\end{description}
\item[(2)] A person shipping a drug under paragraph (a) of this section shall use due diligence to assure that the consignee is regularly engaged in conducting such tests and that the shipment of the new drug will actually be used for tests in vitro or in animals used only for laboratory research.
\item[(3)] A person who ships a drug under paragraph (a) of this section shall maintain adequate records showing the name and post office address of the expert to whom the drug is shipped and the date, quantity, and batch or code mark of each shipment and delivery. Records of shipments under paragraph (a)(1)(i) of this section are to be maintained for a period of 2 years after the shipment. Records and reports of data and shipments under paragraph (a)(1)(ii) of this section are to be maintained in accordance with \S~312.57(b). The person who ships the drug shall upon request from any properly authorized officer or employee of the Food and Drug Administration, at reasonable times, permit such officer or employee to have access to and copy and verify records required to be maintained under this section.
\end{description}

(b) \textit{Termination of authorization to ship.} FDA may terminate authorization to ship a drug under this section if it finds that: (1) the sponsor of the investigation has failed to comply with any of the conditions for shipment established under this section; or (2) the continuance of the investigation is unsafe or otherwise contrary to the public interest or the drug is used for purposes other than bona fide scientific investigation. FDA will notify the person shipping the drug of its finding and invite immediate correction. If correction is not immediately made, the person shall have an opportunity for a regulatory hearing before FDA pursuant to part 16.

(c) \textit{Disposition of unused drug.} The person who ships the drug under paragraph (a) of this section shall assure the return of all unused supplies of the drug from individual investigators whenever the investigation discontinues or the investigation is terminated. The person who ships the drug may authorize in writing alternative disposition of unused supplies of the drug provided this alternative disposition does not expose humans to risks from the drug, either directly or indirectly (e.g., through food-producing animals). The shipper shall maintain records of any alternative disposition.

\subsubsection{Physical AI Adaptation of \S~312.160}

(d) \textit{Physical AI pre-clinical testing.} When investigational drugs are shipped for use in conjunction with Physical AI system testing in laboratory research animals or in vitro, the following additional requirements apply:

\begin{description}[leftmargin=0.8cm,style=sameline]
\item[(i)] The label required under paragraph (a)(1)(i) shall additionally identify any Physical AI system to be used in administering the drug during pre-clinical testing;
\item[(ii)] Records maintained under paragraph (a)(3) shall include documentation of the Physical AI system used in each test, including system model, software version, and the simulation validation status of the system; and
\item[(iii)] Pre-clinical data generated using Physical AI-assisted drug administration shall identify the Physical AI system and its configuration so that the relevance of such data to the planned clinical investigation can be evaluated by FDA.
\end{description}

[52 FR 8831, Mar. 19, 1987, as amended at 52 FR 23031, June 17, 1987. Redesignated at 53 FR 41523, Oct. 21, 1988; 67 FR 9586, Mar. 4, 2002. Physical AI adaptation added 17 March 2026.]

\clearpage

\section{SUBPART H [RESERVED]}

\clearpage

\section{SUBPART I --- EXPANDED ACCESS TO INVESTIGATIONAL DRUGS FOR TREATMENT USE}

\textbf{Source:} 74 FR 40942, Aug. 13, 2009, unless otherwise noted.

\subsection{\S~312.300 General}

(a) \textit{Scope.} This subpart contains the requirements for the use of investigational new drugs and approved drugs where availability is limited by a risk evaluation and mitigation strategy (REMS) when the primary purpose is to diagnose, monitor, or treat a patient's disease or condition. The three categories of expanded access are: (1) expanded access for individual patients, including for emergencies (\S~312.310); (2) expanded access for intermediate-size patient populations (\S~312.315); and (3) expanded access for widespread treatment use through a treatment IND or treatment protocol (\S~312.320).

(b) \textit{Definitions.} The following definitions apply to this subpart:

\begin{description}[leftmargin=0.5cm,style=sameline]
\item[(1)] \textit{Immediately life-threatening disease or condition} means a stage of disease in which there is reasonable likelihood that death will occur within a matter of months or in which premature death is likely without early treatment.
\item[(2)] \textit{Serious disease or condition} means a disease or condition associated with morbidity that has substantial impact on day-to-day functioning. Short-lived and self-limiting morbidity will usually not be sufficient, but the morbidity need not be irreversible if it is persistent or recurrent. Whether a disease or condition is serious is a matter of clinical judgment, based on its impact on such factors as survival, day-to-day functioning, or the likelihood that the disease, if left untreated, will progress from a less severe condition to a more serious one.
\end{description}

(c) \textit{Criteria for expanded access.} FDA must determine that: (1) the patient or patients to be treated have a serious or immediately life-threatening disease or condition, and there is no comparable or satisfactory alternative therapy to diagnose, monitor, or treat the disease or condition; (2) the potential patient benefit justifies the potential risks of the treatment use and those potential risks are not unreasonable in the context of the disease or condition to be treated; and (3) providing the investigational drug for the requested use will not interfere with the initiation, conduct, or completion of clinical investigations that could support marketing approval of the expanded access use or otherwise compromise the potential development of the expanded access use.

(d) \textit{Submission.} Expanded access submissions shall comply with the requirements of \S\S~312.310, 312.315, or 312.320, as applicable.

(e) \textit{Safeguards.} Expanded access to investigational drugs for treatment use is subject to the protections of the IND process, including the requirements for informed consent (part 50) and institutional review board review (part 56), and the IND safety reporting requirements of \S~312.32.

\subsubsection{Physical AI Adaptation of \S~312.300}

(f) \textit{Physical AI expanded access general provisions.} When expanded access to an investigational drug involves the use of a Physical AI system, the following general provisions apply in addition to paragraphs (a) through (e):

\begin{description}[leftmargin=0.8cm,style=sameline]
\item[(i)] The criteria for expanded access under paragraph (c)(2) shall include consideration of the risks and benefits specific to Physical AI-assisted drug administration, including the Physical AI system's safety record from clinical investigations, its current USL score, and the availability of qualified operators at the expanded access site;
\item[(ii)] The informed consent required under paragraph (e) shall include the Physical AI-specific elements described in \S~312.60(f), including a description of the Physical AI system, the human oversight protocols, and the patient's option to request non-Physical AI administration;
\item[(iii)] All Physical AI safety requirements of this part apply to expanded access use, including human oversight (\S~312.60(b)), pre-procedure verification (\S~312.60(c)), emergency stop capabilities, and Physical AI adverse event reporting (\S~312.32(g)); and
\item[(iv)] The sponsor shall maintain the Physical AI system used in expanded access with the same level of maintenance, cybersecurity protection, and software currency as required for clinical investigations under this part.
\end{description}

[74 FR 40942, Aug. 13, 2009. Physical AI adaptation added 17 March 2026.]

\subsection{\S~312.305 Requirements for All Expanded Access Uses}

(a) \textit{General.} The requirements in this section apply to all expanded access uses described in this subpart.

(b) \textit{Submission.} (1) An expanded access submission is required for each type of expanded access described in this subpart. The submission must include: (i) a cover sheet (Form FDA 1571); (ii) the rationale for the intended use of the drug, including a list of available therapeutic options; (iii) sufficient information to establish that the criteria for the specific type of expanded access have been met; (iv) a description of the patient population to be treated; (v) the method of administration of the drug, dose, and duration of therapy; (vi) monitoring procedures; (vii) a description of clinical procedures, laboratory tests, or other monitoring to be used to evaluate the effects of the drug and minimize risk; (viii) chemistry, manufacturing, and controls information; (ix) pharmacology and toxicology information; (x) previous human experience; (xi) additional relevant data from ongoing or other trials; and (xii) a commitment to provide a written report within a specified time following completion or discontinuation.

\begin{description}[leftmargin=0.5cm,style=sameline]
\item[(2)] The expanded access submission and its contents constitute the IND for purposes of this part, unless the expanded access use is added to an existing IND, in which case the existing IND number shall be referenced.
\end{description}

(c) \textit{Safeguards.} (1) Treatment is generally overseen by an investigator or treating physician who is qualified by training and experience. (2) An IRB must review and approve the expanded access use, except for individual patient access under \S~312.310 when the treating physician determines that the need for treatment outweighs the need for IRB review, with written documentation of the determination. (3) Informed consent requirements of part 50 apply. (4) IND safety reporting requirements of \S~312.32 apply. (5) FDA may require any sponsor of an expanded access IND or protocol to monitor the use if necessary.

(d) \textit{Charging for expanded access.} A sponsor or investigator may charge for an investigational drug under expanded access as described in \S~312.8.

\subsubsection{Physical AI Adaptation of \S~312.305}

(e) \textit{Physical AI expanded access submissions.} When an expanded access submission involves the use of a Physical AI system, the submission required under paragraph (b)(1) shall additionally include:

\begin{description}[leftmargin=0.8cm,style=sameline]
\item[(i)] The Physical AI System Description or a summary thereof, or a reference to a Physical AI System Description already on file with an existing IND;
\item[(ii)] The current USL score for each Physical AI system to be used in the expanded access;
\item[(iii)] A description of the Physical AI system configuration at each expanded access site, including hardware model, software version, and AI/ML model version;
\item[(iv)] Documentation of the qualifications and training of the physicians who will operate or supervise the Physical AI system at each expanded access site; and
\item[(v)] A description of the Physical AI monitoring procedures, including the frequency and method of Physical AI performance review, Physical AI adverse event reporting procedures, and the criteria for discontinuing Physical AI-assisted treatment.
\end{description}

(f) \textit{Physical AI expanded access safeguards.} In addition to the safeguards in paragraph (c):

\begin{description}[leftmargin=0.8cm,style=sameline]
\item[(i)] IRB review of expanded access involving a Physical AI system shall include review of the Physical AI system's role, its safety record, and the informed consent language regarding the Physical AI system;
\item[(ii)] Physical AI adverse event reporting under \S~312.32(g) applies to all expanded access uses involving Physical AI systems; and
\item[(iii)] The sponsor shall designate a qualified individual to provide remote Physical AI system technical support to expanded access sites during Physical AI-assisted procedures, unless the site has on-site Physical AI system engineering support.
\end{description}

[74 FR 40942, Aug. 13, 2009. Physical AI adaptation added 17 March 2026.]

\subsection{\S~312.310 Individual Patient Access to Investigational Drugs for Treatment Use, Including for Emergencies}

(a) \textit{Criteria.} Under this section, FDA may permit an investigational drug to be used for the treatment of an individual patient by a licensed physician. FDA must determine that:

\begin{description}[leftmargin=0.5cm,style=sameline]
\item[(1)] The patient has a serious or immediately life-threatening disease or condition, and there is no comparable or satisfactory alternative drug or other therapy available to treat that stage of the disease or condition in the patient;
\item[(2)] The potential patient benefit justifies the potential risks of the treatment use and those potential risks are not unreasonable in the context of the disease or condition to be treated; and
\item[(3)] Providing the investigational drug for the requested use will not interfere with the initiation, conduct, or completion of clinical investigations that could support marketing approval of the expanded access use or otherwise compromise the potential development of the expanded access use.
\end{description}

(b) \textit{Submission.} The request for individual patient access must be submitted by a licensed physician who will be the treating physician. The request may be made by telephone or other rapid means of communication, with a written submission to follow within 15 working days. The request shall include: (1) a description of the patient's disease or condition; (2) a summary of prior therapy; (3) a brief clinical protocol for the proposed use; (4) the method of administration of the drug, dose, and duration of therapy; (5) a statement of the physician's qualifications to use the investigational drug for the proposed treatment; (6) patient informed consent document; and (7) evidence that the drug manufacturer or sponsor is aware of the proposed expanded access use.

(c) \textit{Safeguards.} Treatment is generally limited to a single course of therapy for a specified duration, unless FDA authorizes multiple courses. The physician must report adverse events within 15 calendar days.

(d) \textit{Emergency procedures.} If there is an emergency that requires the patient to be treated before a written submission can be made, FDA may authorize treatment by telephone. The physician shall submit an expanded access submission within 15 working days.

\subsubsection{Physical AI Adaptation of \S~312.310}

(e) \textit{Physical AI in individual patient access.} When an investigational drug is to be administered through or with the assistance of a Physical AI system under individual patient access:

\begin{description}[leftmargin=0.8cm,style=sameline]
\item[(i)] The submission shall include a brief description of the Physical AI system to be used, its current USL score, and the physician's training and experience with the specific Physical AI system;
\item[(ii)] The patient benefit-risk determination under paragraph (a)(2) shall consider the risks and benefits of Physical AI-assisted administration compared to manual administration, including whether the Physical AI system provides clinically meaningful advantages in precision, consistency, or safety for the specific patient's condition;
\item[(iii)] The informed consent document shall describe the Physical AI system's role in the treatment, the human oversight protocols, and the patient's option to request non-Physical AI administration if clinically feasible; and
\item[(iv)] In emergency situations under paragraph (d), the physician may proceed with Physical AI-assisted treatment under the telephone authorization, provided the Physical AI system has a valid USL score above the applicable threshold and the physician has current training on the system.
\end{description}

[74 FR 40942, Aug. 13, 2009. Physical AI adaptation added 17 March 2026.]

\subsection{\S~312.315 Intermediate-Size Patient Populations}

(a) \textit{Criteria.} Under this section, FDA may permit an investigational drug to be used for the treatment of a patient population smaller than that typical of a treatment IND or treatment protocol. FDA must determine that:

\begin{description}[leftmargin=0.5cm,style=sameline]
\item[(1)] There is enough evidence that the drug is safe at the dose and duration proposed for expanded access use to justify a clinical trial of the drug in the approximate number of patients expected to be treated;
\item[(2)] There is at least preliminary clinical evidence of effectiveness of the drug, or of a plausible pharmacologic effect of the drug to make expanded access use a reasonable therapeutic option in the anticipated patient population; and
\item[(3)] The criteria in \S~312.310(a) are met.
\end{description}

(b) \textit{Submission.} The expanded access submission shall include information adequate to satisfy the criteria in paragraph (a), including: (1) a description of the patient population; (2) the method of administration, dose, and duration; (3) monitoring procedures; (4) a description of available safety data; (5) clinical protocol; and (6) investigator information.

(c) \textit{Safeguards.} The sponsor shall monitor the expanded access use, ensure adverse events are reported, and ensure investigators are informed of new safety information.

\subsubsection{Physical AI Adaptation of \S~312.315}

(d) \textit{Physical AI in intermediate-size access.} When a Physical AI system is used to administer an investigational drug under intermediate-size patient population access:

\begin{description}[leftmargin=0.8cm,style=sameline]
\item[(i)] The expanded access submission shall include the Physical AI System Description or a summary sufficient to evaluate the Physical AI system's safety for the proposed use, the current USL score, and a description of the sites where the Physical AI system will be used and the qualifications of the treating physicians;
\item[(ii)] The monitoring requirements under paragraph (c) shall include Physical AI system performance monitoring and Physical AI adverse event reporting consistent with the requirements of \S~312.32(g);
\item[(iii)] All Physical AI safety requirements of this part apply to expanded access use, including human oversight, pre-procedure verification, and emergency stop requirements; and
\item[(iv)] The sponsor shall collect Physical AI performance data from the expanded access use and make such data available for inclusion in any subsequent marketing application.
\end{description}

[74 FR 40942, Aug. 13, 2009. Physical AI adaptation added 17 March 2026.]

\subsection{\S~312.320 Treatment IND or Treatment Protocol}

(a) \textit{Criteria.} Under this section, FDA may permit an investigational drug to be used for widespread treatment use. FDA must determine that:

\begin{description}[leftmargin=0.5cm,style=sameline]
\item[(1)] The drug is being investigated in a controlled clinical trial under an IND designed to support a marketing application for the expanded access use, or all clinical trials of the drug have been completed;
\item[(2)] The sponsor of the controlled clinical trial is actively pursuing marketing approval with due diligence;
\item[(3)] There is sufficient clinical evidence of safety and effectiveness to support the expanded access use; and
\item[(4)] The criteria in \S~312.310(a) are met.
\end{description}

(b) \textit{Submission.} The expanded access submission shall include information adequate to satisfy the criteria in paragraph (a), including a clinical protocol for the treatment use, an investigator's brochure, safety data summaries, FDA status of the drug, and a statement regarding marketing application progress.

(c) \textit{Safeguards.} The sponsor shall monitor the treatment protocol, comply with IND safety reporting requirements, ensure investigators are informed of safety data, and submit annual reports.

\subsubsection{Physical AI Adaptation of \S~312.320}

(d) \textit{Physical AI in treatment INDs.} When a Physical AI system is used to administer an investigational drug under a treatment IND or treatment protocol:

\begin{description}[leftmargin=0.8cm,style=sameline]
\item[(i)] The treatment protocol shall include a Physical AI System Description or reference to a Physical AI System Description already on file with the IND, and shall describe the Physical AI system configuration to be used at treatment sites;
\item[(ii)] The expanded access submission shall include evidence that the Physical AI system has performed safely in clinical investigations under the IND, including aggregate Physical AI adverse event data and USL score trends from completed or ongoing trials;
\item[(iii)] Treatment sites shall meet the Physical AI readiness requirements of \S~312.40(e), including investigator training, pre-procedure safety matrix verification, and MCP infrastructure configuration;
\item[(iv)] The sponsor's monitoring requirements under paragraph (c) shall include cross-site Physical AI performance monitoring to detect any site-specific Physical AI performance issues;
\item[(v)] Physical AI adverse events occurring during treatment use shall be reported in accordance with \S~312.32(g) and shall be included in the annual report required under paragraph (c); and
\item[(vi)] The sponsor shall maintain the Physical AI system at treatment sites with the same level of cybersecurity protection, software updates, and maintenance as required for the controlled clinical trial.
\end{description}

[74 FR 40942, Aug. 13, 2009. Physical AI adaptation added 17 March 2026.]

\clearpage

\section{SUBPART J --- PHYSICAL AI SYSTEMS IN CLINICAL INVESTIGATIONS (NEW)}

\textit{Note: This subpart is entirely new and has no corresponding subpart in the original 21 CFR Part 312. It consolidates and expands upon the Physical AI-specific requirements introduced throughout this adapted regulation.}

\subsection{\S~312.400 Purpose and Scope of Subpart J}

(a) \textit{Purpose.} This subpart establishes the comprehensive regulatory framework for the use of Physical AI systems in clinical investigations of investigational new drugs. It consolidates the Physical AI-specific requirements referenced throughout this part and provides additional detailed requirements for Physical AI system validation, operation, monitoring, and lifecycle management during clinical investigations.

(b) \textit{Scope.} This subpart applies to all clinical investigations conducted under this part in which a Physical AI system, as defined in \S~312.1(b), is used to assist in, perform, or monitor any aspect of investigational drug administration, including but not limited to:

\begin{description}[leftmargin=0.8cm,style=sameline]
\item[(1)] Robotic surgical systems that administer investigational drugs during surgical procedures;
\item[(2)] Automated drug delivery systems that prepare, measure, or deliver investigational drugs;
\item[(3)] AI-guided treatment planning systems that use digital twin models, simulation data, or machine learning algorithms to determine drug administration parameters;
\item[(4)] Robotic monitoring systems that assess patient response to investigational drugs through automated vital sign measurement, biomarker sampling, or imaging; and
\item[(5)] Integrated Physical AI platforms that combine two or more of the above functions.
\end{description}

(c) \textit{Relationship to other subparts.} The requirements of this subpart supplement, and do not replace, the requirements of Subparts A through I. In the event of a conflict between a provision of this subpart and a Physical AI adaptation provision in another subpart, the more specific provision shall control.

[Physical AI adaptation added 17 March 2026.]

\subsection{\S~312.401 Physical AI System Classification for Clinical Investigations}

(a) \textit{Risk-based classification.} Physical AI systems used in clinical investigations shall be classified according to their level of autonomy and clinical risk:

\begin{description}[leftmargin=0.8cm,style=sameline]
\item[(1)] \textit{Class I --- Assistive Physical AI.} Systems that provide decision support, guidance, or enhanced visualization to the investigator but do not directly manipulate the investigational drug or the patient. The investigator retains full manual control of drug administration at all times. Examples include AI-guided imaging systems that assist in identifying optimal injection sites and digital twin systems used solely for treatment planning without automated execution.
\item[(2)] \textit{Class II --- Collaborative Physical AI.} Systems that perform specific drug administration tasks under continuous direct supervision by a qualified investigator who maintains the ability to override or halt the system's actions in real time. The system and investigator share control of the drug administration process. Examples include robotic infusion systems that execute physician-approved infusion protocols and robotic surgical systems that perform drug delivery under surgeon guidance.
\item[(3)] \textit{Class III --- Supervised Autonomous Physical AI.} Systems that can execute complete drug administration procedures with limited investigator intervention, under supervisory oversight. The investigator monitors the system's performance and intervenes when specified conditions are met, but does not control each individual action. Examples include automated closed-loop drug delivery systems that adjust dosing based on real-time patient monitoring data and robotic brachytherapy systems that autonomously place radioactive sources according to a pre-approved treatment plan.
\end{description}

(b) \textit{Classification requirements.} The sponsor shall classify each Physical AI system in the Physical AI System Description submitted under \S~312.23(g). FDA may reclassify a Physical AI system if it determines that the sponsor's classification does not accurately reflect the system's level of autonomy or clinical risk.

(c) \textit{Class-specific requirements.}

\begin{description}[leftmargin=0.8cm,style=sameline]
\item[(1)] \textit{Class I systems} are subject to the general Physical AI requirements of this part, including the Physical AI System Description, simulation validation, and Physical AI adverse event reporting requirements.
\item[(2)] \textit{Class II systems} are subject to all Class I requirements and, in addition, must demonstrate USL scores meeting or exceeding the thresholds specified in \S~312.1(e), must implement real-time safety monitoring with automatic safety stop capabilities, and must provide the investigator with real-time feedback on system actions and drug administration status.
\item[(3)] \textit{Class III systems} are subject to all Class I and Class II requirements and, in addition, must demonstrate enhanced simulation validation including multi-scenario failure mode testing, must implement redundant safety systems with independent monitoring channels, must maintain continuous telemetry logging at a resolution sufficient to reconstruct all system decisions and actions, and must undergo periodic re-validation during the clinical investigation as specified in the PCCP.
\end{description}

[Physical AI adaptation added 17 March 2026.]

\subsection{\S~312.402 Physical AI System Validation Requirements}

(a) \textit{Pre-clinical validation.} Before a Physical AI system may be used in a clinical investigation, the sponsor shall complete the following validation activities:

\begin{description}[leftmargin=0.8cm,style=sameline]
\item[(1)] \textit{Simulation validation.} The Physical AI system shall be validated through simulation testing as described in \S~312.21(d), including validation across multiple simulation frameworks to demonstrate cross-framework behavioral consistency. The simulation validation shall test the system across the full range of clinical scenarios anticipated in the protocol, including normal operating conditions, boundary conditions, and failure scenarios;
\item[(2)] \textit{Bench testing.} The Physical AI system shall undergo bench testing using physical phantoms, tissue analogues, or other non-clinical test articles that replicate the mechanical, thermal, and chemical properties of the target clinical environment. Bench testing shall include verification of drug delivery accuracy, force and torque limits, sensor calibration, and emergency stop response times;
\item[(3)] \textit{Integration testing.} The Physical AI system shall be tested in its complete clinical configuration, including all hardware components, software modules, network connections, and peripheral devices, to verify that all system components function correctly in concert and that no integration defects compromise safety or performance; and
\item[(4)] \textit{Sim-to-real gap assessment.} The sponsor shall quantify the sim-to-real gap by comparing the Physical AI system's performance in simulation to its performance in bench testing, identifying any discrepancies, and establishing tolerance ranges for acceptable divergence. The sim-to-real gap assessment shall be updated as the system is modified during the investigation.
\end{description}

(b) \textit{Clinical site validation.} Before a Physical AI system may be used in clinical procedures at a specific investigational site, the following site-specific validation shall be completed:

\begin{description}[leftmargin=0.8cm,style=sameline]
\item[(1)] \textit{Installation qualification.} Verification that the Physical AI system has been installed at the site in accordance with the manufacturer's specifications and the Physical AI System Description, including environmental conditions (temperature, humidity, vibration, electromagnetic interference), power supply and backup power systems, network connectivity and bandwidth, and physical space and access requirements;
\item[(2)] \textit{Operational qualification.} Verification that the Physical AI system operates within the specified performance parameters at the site, including execution of the standard test procedures defined in the Physical AI System Description, verification of all safety systems including emergency stop, force limits, and collision detection, and confirmation of communication with all required peripheral devices and monitoring systems; and
\item[(3)] \textit{Performance qualification.} Demonstration that the Physical AI system can consistently perform the clinical procedures authorized under the protocol, using site-specific test articles, within the acceptance criteria defined in the Physical AI System Description.
\end{description}

(c) \textit{Ongoing validation.} During the clinical investigation, the sponsor shall:

\begin{description}[leftmargin=0.8cm,style=sameline]
\item[(1)] Monitor Physical AI system performance against the pre-specified acceptance criteria established during pre-clinical validation;
\item[(2)] Re-validate the Physical AI system following any software update, AI/ML model update, hardware modification, or safety event, in accordance with the PCCP and the procedures described in the Physical AI System Description; and
\item[(3)] Conduct periodic re-validation at intervals specified in the protocol, including USL reassessment and simulation re-validation, to confirm that the system's performance has not degraded over time.
\end{description}

[Physical AI adaptation added 17 March 2026.]

\subsection{\S~312.403 Physical AI System Cybersecurity Requirements}

(a) \textit{Cybersecurity by design.} Physical AI systems used in clinical investigations shall incorporate cybersecurity protections in their design, including:

\begin{description}[leftmargin=0.8cm,style=sameline]
\item[(1)] \textit{Authentication and access control.} Multi-factor authentication for operator access, role-based access controls that limit system functions to authorized personnel, and automatic session timeout after periods of inactivity;
\item[(2)] \textit{Data protection.} Encryption of data in transit and at rest, including telemetry data, patient data, drug administration records, and system configuration data. Encryption shall use algorithms and key lengths consistent with current NIST guidelines;
\item[(3)] \textit{Network security.} Network segmentation isolating the Physical AI system from general hospital networks, intrusion detection and prevention systems monitoring Physical AI network traffic, and firewall configurations restricting inbound and outbound connections to authorized endpoints;
\item[(4)] \textit{Software integrity.} Secure boot processes that verify firmware and software integrity at startup, code signing for all software components with verification before execution, and integrity monitoring that detects unauthorized modifications to system software; and
\item[(5)] \textit{Audit trails.} Hash-chained, cryptographically signed audit trails that record all system events, operator actions, configuration changes, and network communications in a manner that prevents undetected tampering.
\end{description}

(b) \textit{Vulnerability management.} The sponsor shall maintain a vulnerability management program for the Physical AI system, including:

\begin{description}[leftmargin=0.8cm,style=sameline]
\item[(1)] Regular vulnerability scanning and assessment;
\item[(2)] A defined process for evaluating and applying security patches, with timelines proportionate to the severity of the vulnerability;
\item[(3)] A software bill of materials (SBOM) identifying all third-party software components and their versions; and
\item[(4)] Monitoring of vulnerability databases and threat intelligence sources for vulnerabilities affecting Physical AI system components.
\end{description}

(c) \textit{Incident response.} The sponsor shall maintain a cybersecurity incident response plan that includes:

\begin{description}[leftmargin=0.8cm,style=sameline]
\item[(1)] Procedures for detecting, containing, and remediating cybersecurity incidents affecting the Physical AI system;
\item[(2)] Criteria for determining when a cybersecurity incident constitutes a Physical AI adverse event reportable under \S~312.32(g);
\item[(3)] Communication procedures for notifying affected investigational sites, FDA, and other relevant parties; and
\item[(4)] Procedures for preserving forensic evidence from cybersecurity incidents for analysis and regulatory review.
\end{description}

[Physical AI adaptation added 17 March 2026.]

\subsection{\S~312.404 Physical AI Human Oversight Requirements}

(a) \textit{General oversight requirement.} All Physical AI-assisted clinical procedures involving investigational drugs shall be conducted under human oversight by a qualified investigator or subinvestigator. No Physical AI system shall administer an investigational drug to a human subject without human oversight, regardless of the system's classification under \S~312.401.

(b) \textit{Levels of oversight.} The level of human oversight required shall correspond to the Physical AI system's classification:

\begin{description}[leftmargin=0.8cm,style=sameline]
\item[(1)] \textit{Class I systems.} The investigator maintains full manual control and uses the Physical AI system's output as advisory information. No additional oversight requirements beyond standard clinical practice.
\item[(2)] \textit{Class II systems.} The investigator maintains continuous direct supervision of the Physical AI system during all drug administration activities. The investigator must be physically present in the procedure room, must have immediate access to manual override controls, and must be able to halt the system and assume manual control within a response time specified in the protocol (not to exceed 2 seconds for drug delivery systems and 500 milliseconds for surgical systems).
\item[(3)] \textit{Class III systems.} In addition to the Class II requirements, a second qualified individual shall be present during all drug administration procedures to serve as an independent safety monitor. The safety monitor's sole responsibility during the procedure shall be to observe the Physical AI system's performance and alert the supervising investigator to any anomalies. The protocol shall define the specific performance indicators the safety monitor is required to observe.
\end{description}

(c) \textit{Operator-to-system ratio.} No investigator or subinvestigator shall simultaneously supervise more than one Physical AI system performing drug administration procedures. The protocol may specify a higher operator-to-system ratio (i.e., more operators per system) based on procedure complexity or risk.

(d) \textit{Fatigue management.} The protocol shall include provisions to prevent operator fatigue from compromising human oversight effectiveness, including maximum continuous supervision periods and mandatory rest intervals for Physical AI system operators during extended procedures.

(e) \textit{Override and emergency stop.}

\begin{description}[leftmargin=0.8cm,style=sameline]
\item[(1)] Every Physical AI system used in clinical investigations shall be equipped with a physical emergency stop mechanism that is immediately accessible to the supervising investigator and that, when activated, brings the system to a safe state within a time period appropriate for the clinical procedure (not to exceed 500 milliseconds for all moving components);
\item[(2)] The emergency stop mechanism shall function independently of the Physical AI system's software and shall not be subject to software override or delay;
\item[(3)] The Physical AI system shall provide the supervising investigator with a manual override capability that allows the investigator to assume direct control of any system function at any time during a procedure; and
\item[(4)] The protocol shall define the criteria under which the investigator shall activate the emergency stop or manual override, including but not limited to: unexpected patient physiological response, Physical AI system behavior outside pre-specified parameters, loss of system sensor data, and loss of communication with the Physical AI system.
\end{description}

[Physical AI adaptation added 17 March 2026.]

\subsection{\S~312.405 Physical AI System Lifecycle Management}

(a) \textit{Configuration management.} The sponsor shall maintain configuration management records for each Physical AI system used in the investigation, documenting:

\begin{description}[leftmargin=0.8cm,style=sameline]
\item[(1)] The complete hardware configuration, including all components, serial numbers, firmware versions, and calibration dates;
\item[(2)] The complete software configuration, including operating system version, application software versions, AI/ML model versions and training data provenance, and all configuration parameters; and
\item[(3)] All changes to the hardware or software configuration, including the date of the change, the nature of the change, the rationale, the validation performed, and any required regulatory submissions under \S~312.30.
\end{description}

(b) \textit{AI/ML model management.} When the Physical AI system includes AI/ML models that influence drug administration decisions or actions:

\begin{description}[leftmargin=0.8cm,style=sameline]
\item[(1)] The sponsor shall document the training data, training methodology, performance metrics, and validation results for each AI/ML model;
\item[(2)] Model updates shall be managed in accordance with the PCCP submitted under \S~312.23(g)(vi), including pre-specified performance criteria for accepting or rejecting model updates;
\item[(3)] The sponsor shall monitor for model drift, data drift, and concept drift that could affect model performance in the clinical setting, and shall implement corrective actions when drift exceeds pre-specified thresholds; and
\item[(4)] The sponsor shall maintain a complete model version history, including the ability to revert to any previous model version if a new model version demonstrates unacceptable performance.
\end{description}

(c) \textit{Decommissioning.} When a Physical AI system is removed from a clinical investigation, the sponsor shall:

\begin{description}[leftmargin=0.8cm,style=sameline]
\item[(1)] Complete all required data archival, including transfer of telemetry data, system logs, and audit trails to the sponsor's long-term storage;
\item[(2)] Securely erase all patient data, investigational drug data, and access credentials from the Physical AI system;
\item[(3)] Revoke all network access, authentication tokens, and MCP server configurations associated with the system;
\item[(4)] Document the decommissioning in the Physical AI system maintenance record; and
\item[(5)] Retain the Physical AI system or key components for the record retention period specified in \S~312.57(f), or document the disposition of the system if it is returned to the manufacturer or otherwise disposed of.
\end{description}

[Physical AI adaptation added 17 March 2026.]

\clearpage

\section*{REFERENCES AND BIBLIOGRAPHY}

\addcontentsline{toc}{section}{References and Bibliography}

The following references informed the development of this Physical AI adaptation of 21 CFR Part 312. This bibliography includes the original regulatory sources, FDA guidance documents, relevant standards, and the Physical AI and oncology literature that provides the technical foundation for the adapted provisions.

\subsection*{Primary Regulatory Sources}

\begin{enumerate}
\item 21 CFR Part 312 --- Investigational New Drug Application. Code of Federal Regulations, Title 21, Chapter I, Subchapter D, Part 312. U.S. Food and Drug Administration.

\item Federal Food, Drug, and Cosmetic Act (FD\&C Act), 21 U.S.C. \S\S~301--399i, as amended.

\item Public Health Service Act, 42 U.S.C. \S~262 (Section 351), as amended.

\item 21 CFR Part 50 --- Protection of Human Subjects. Code of Federal Regulations, Title 21, Chapter I, Subchapter A.

\item 21 CFR Part 56 --- Institutional Review Boards. Code of Federal Regulations, Title 21, Chapter I, Subchapter A.

\item 21 CFR Part 11 --- Electronic Records; Electronic Signatures. Code of Federal Regulations, Title 21, Chapter I, Subchapter A.
\end{enumerate}

\subsection*{FDA Guidance Documents}

\begin{enumerate}[resume]
\item U.S. Food and Drug Administration. ``Artificial Intelligence and Machine Learning (AI/ML)-Enabled Device Software Functions: Lifecycle Management and Marketing Submission Recommendations.'' Draft Guidance, January 2025.

\item U.S. Food and Drug Administration. ``Marketing Submission Recommendations for a Predetermined Change Control Plan for Artificial Intelligence/Machine Learning (AI/ML)-Enabled Device Software Functions.'' Guidance for Industry, December 2024.

\item U.S. Food and Drug Administration. ``Cybersecurity in Medical Devices: Quality System Considerations and Content of Premarket Submissions.'' Guidance for Industry, September 2023.

\item U.S. Food and Drug Administration. ``Computer Software Assurance for Production and Quality System Software.'' Guidance for Industry, September 2022.

\item U.S. Food and Drug Administration. ``Clinical Decision Support Software.'' Guidance for Industry and FDA Staff, September 2022.

\item U.S. Food and Drug Administration. ``Content of Premarket Submissions for Device Software Functions.'' Guidance for Industry, June 2023.

\item U.S. Food and Drug Administration. ``Investigational Device Exemptions (IDEs) for Early Feasibility Medical Device Clinical Studies, Including Certain First in Human (FIH) Studies.'' Guidance for Industry, October 2013.
\end{enumerate}

\subsection*{Robotics and Physical AI Standards}

\begin{enumerate}[resume]
\item ISO 13482:2014. Robots and robotic devices --- Safety requirements for personal care robots. International Organization for Standardization.

\item ISO 10218-1:2011. Robots and robotic devices --- Safety requirements for industrial robots --- Part 1: Robots. International Organization for Standardization.

\item IEC 62304:2006/AMD 1:2015. Medical device software --- Software life cycle processes. International Electrotechnical Commission.

\item IEC 80001-1:2021. Application of risk management for IT-networks incorporating medical devices. International Electrotechnical Commission.

\item IEEE 7001-2021. IEEE Standard for Transparency of Autonomous Systems. Institute of Electrical and Electronics Engineers.

\item IEEE 7010-2020. IEEE Recommended Practice for Assessing the Impact of Autonomous and Intelligent Systems on Human Well-Being. Institute of Electrical and Electronics Engineers.

\item NIST AI 100-1. Artificial Intelligence Risk Management Framework (AI RMF 1.0). National Institute of Standards and Technology, January 2023.
\end{enumerate}

\subsection*{Physical AI and Simulation Literature}

\begin{enumerate}[resume]
\item Kawchak, K. ``Physical AI Oncology Trials: A Framework for Integrating Robotic Systems into Clinical Drug Investigations.'' GitHub Repository, 2025--2026. Available at: https://github.com/kevinkawchak/physical-ai-oncology-trials.

\item NVIDIA Corporation. ``Isaac Sim: Robotics Simulation and Synthetic Data Generation.'' NVIDIA Developer Documentation, 2024--2025.

\item NVIDIA Corporation. ``NVIDIA Cosmos: World Foundation Models for Physical AI.'' Technical Report, 2025.

\item Google DeepMind. ``Gemini Robotics: Bringing AI into the Physical World.'' Technical Report, March 2025.

\item Todorov, E., Erez, T., and Tassa, Y. ``MuJoCo: A physics engine for model-based control.'' In Proceedings of IEEE/RSJ International Conference on Intelligent Robots and Systems, 2012.

\item Makoviychuk, V., et al. ``Isaac Gym: High Performance GPU-Based Physics Simulation for Robot Learning.'' In Proceedings of the Neural Information Processing Systems (NeurIPS) Datasets and Benchmarks Track, 2021.

\item Tobin, J., et al. ``Domain Randomization for Transferring Deep Neural Networks from Simulation to the Real World.'' In Proceedings of IEEE/RSJ International Conference on Intelligent Robots and Systems, 2017.

\item Zhao, W., et al. ``Sim-to-Real Transfer in Deep Reinforcement Learning for Robotics: A Survey.'' IEEE Transactions on Neural Networks and Learning Systems, 2020.
\end{enumerate}

\subsection*{Oncology Robotics and Drug Delivery}

\begin{enumerate}[resume]
\item Trejos, A.L., et al. ``Robot-assisted minimally invasive surgery and interventional oncology.'' In Handbook of Robotic and Image-Guided Surgery, Elsevier, 2020.

\item Fichtinger, G., et al. ``Robotic assistance for ultrasound-guided brachytherapy.'' Medical Image Analysis, 2008.

\item Podder, T.K., et al. ``AAPM and GEC-ESTRO guidelines for image-guided robotic brachytherapy: Report of Task Group 192.'' Medical Physics, 2014.

\item Yang, G.-Z., et al. ``Medical robotics --- Regulatory, ethical, and legal considerations for increasing levels of autonomy.'' Science Robotics, 2017.

\item Dupont, P.E., et al. ``A decade retrospective of medical robotics research from 2010 to 2020.'' Science Robotics, 2021.
\end{enumerate}

\subsection*{AI/ML in Clinical Trials}

\begin{enumerate}[resume]
\item Topol, E.J. ``High-performance medicine: the convergence of human and artificial intelligence.'' Nature Medicine, 25(1):44--56, 2019.

\item Liu, X., et al. ``Reporting guidelines for clinical trial reports for interventions involving artificial intelligence: the CONSORT-AI extension.'' The Lancet Digital Health, 2(10):e537--e548, 2020.

\item Cruz Rivera, S., et al. ``Guidelines for clinical trial protocols for interventions involving artificial intelligence: the SPIRIT-AI extension.'' The Lancet Digital Health, 2(10):e549--e560, 2020.

\item U.S. Food and Drug Administration. ``Using Artificial Intelligence and Machine Learning in the Development of Drug and Biological Products.'' Discussion Paper, May 2023.

\item Anthropic. ``Model Card and Evaluations for Claude Models.'' Technical Documentation, 2024--2025.
\end{enumerate}

\subsection*{Cybersecurity in Medical Systems}

\begin{enumerate}[resume]
\item NIST Special Publication 800-53, Rev. 5. ``Security and Privacy Controls for Information Systems and Organizations.'' National Institute of Standards and Technology, September 2020.

\item NIST Special Publication 800-183. ``Networks of `Things'.'' National Institute of Standards and Technology, July 2016.

\item Health Care Industry Cybersecurity Task Force. ``Report on Improving Cybersecurity in the Health Care Industry.'' U.S. Department of Health and Human Services, June 2017.

\item Medical Device Innovation Consortium (MDIC). ``Medical Device Cybersecurity: Regional Incident Preparedness and Response Playbook.'' 2018.
\end{enumerate}

\subsection*{Digital Twins and Treatment Planning}

\begin{enumerate}[resume]
\item Lal, A., et al. ``Digital twins and AI in clinical trials.'' Nature Medicine, 30:3405--3407, 2024.

\item Corral-Acero, J., et al. ``The `Digital Twin' to enable the vision of precision cardiology.'' European Heart Journal, 41(48):4556--4564, 2020.

\item Hernandez-Boussard, T., et al. ``Digital twins for predictive oncology will be a paradigm shift for precision cancer care.'' Nature Medicine, 27:2065--2066, 2021.

\item Venkatesh, K.P., et al. ``Health digital twins as tools for precision medicine: Considerations for computation, implementation, and regulation.'' npj Digital Medicine, 5:150, 2022.
\end{enumerate}

\vspace{1cm}

\noindent\rule{\textwidth}{0.4pt}

\vspace{0.5cm}

\begin{center}
\textbf{END OF ADAPTED REGULATION}

\vspace{0.3cm}

\textit{This adaptation of 21 CFR Part 312 for Physical AI systems in oncology clinical trials was prepared as part of the Physical AI Oncology Trials project.}

\vspace{0.2cm}

\textit{Version 1.0 --- 17 March 2026}

\vspace{0.2cm}

\textit{This document is intended for research and educational purposes. It does not represent official FDA regulation or guidance.}
\end{center}

\end{document}
